\documentclass{beamer}					% Document class

\usepackage[english]{babel}				% Set language
\usepackage[utf8x]{inputenc}			% Set encoding

\mode<presentation>						% Set options
{
  \usetheme{default}					% Set theme
  \usecolortheme{default} 				% Set colors
  \usefonttheme{default}  				% Set font theme
  \setbeamertemplate{caption}[numbered]	% Set caption to be numbered
}

% Uncomment this to have the outline at the beginning of each section highlighted.
%\AtBeginSection[]
%{
%  \begin{frame}{Outline}
%    \tableofcontents[currentsection]
%  \end{frame}
%}

\usepackage{graphicx}					% For including figures
\usepackage{booktabs}					% For table rules
\usepackage{hyperref}					% For cross-referencing

\title{DM8100 - High Performance Computing}	    % Presentation title
\author{Ask Madsen, Bjarke Hammerbak Paluszewski}								% Presentation author
\institute{University of Southern Denmark}					% Author affiliation
\date{\today}									% Today's date	
\begin{document}

% Title page
% This page includes the informations defined earlier including title, author/s, affiliation/s and the date
\begin{frame}
  \titlepage
\end{frame}

% Outline
% This page includes the outline (Table of content) of the presentation. All sections and subsections will appear in the outline by default.
\begin{frame}{Outline}
  \tableofcontents
\end{frame}

% The following is the most frequently used slide types in beamer
% The slide structure is as follows:
%
%\begin{frame}{<slide-title>}
%	<content>
%\end{frame}

\section{Section One}

\begin{frame}{Slide with bullet points}
	This is a bullet list of two points:
    \begin{itemize}
		\item Point one
        \item Point two
	\end{itemize}
\end{frame}

\begin{frame}{Slide with two columns}
	\begin{columns}
		\column{.5\textwidth}
        Text goes in first column.
        
        \column{.5\textwidth}
        Text goes in second column
	\end{columns}
\end{frame}

\section{Section Two}

\end{document}